\section{Discussion and Conclusion}
\label{sec:conclusion}

We presented a systematic adversarial audit of approximate graph
unlearning. Deletion-set selection alone exposes a 2D
\emph{Vulnerability Fingerprint}: GNNDelete is highly vulnerable, GIF is
weakly TracIn-aligned, the Partition pair exhibits Shard Protection, and
MEGU/IDEA remain near the origin despite belonging to vulnerable parent
families. MEGU and IDEA should therefore be read as method-specific nulls
under the tested IM/TracIn selector class, not as a new robust family;
mechanism-aware selectors for mutual-evolution and certified-gradient
objectives are the natural next test.

\paragraph{Design lessons.}
Shard Protection suggests that structural isolation can act as a
firewall against cascading approximation error, while TracIn-sensitive
methods show the risk of unconstrained gradient alignment. Future GU
systems should report adversarial-deletion bounds alongside runtime and
random-deletion utility, and should explore soft partition barriers or
minimax-robust update objectives.

\paragraph{Limitations.}
\label{sec:limitations}
Our shard-count sweep is partial (\S\ref{app:shard}); ogbn-arxiv uses
three seeds; TracIn's $G$-matrix does not scale to very high feature
dimension; ratio sweeps are small-graph only; and we cover three of four
OpenGU categories. The L2-\textsc{direct} axis is an upper bound on
shadow-model L2-\textsc{surrogate} attacks, motivated by standard
transferability results~\cite{papernot2017practical,demontis2019adversarial}
but not yet quantified across graph architectures. Because the full
design spans seven axes, we use claim-conditioned slices rather than a
Cartesian-complete grid (Appendix~\ref{app:results}). Finally, our
update-detection AUC is a posterior-shift deletion-membership audit, not
a full shadow-model MIA.

\paragraph{Broader impact and future work.}
\label{sec:broader}
\label{sec:future}
The diagnostic kit here ($\Delta F$ decomposition, retrain gap,
prediction shift, hop-distance decay, and update-detection AUC) can be
reused with stronger selectors. A coalition can reach L1 structural
attacks at zero model-access cost, while L2-\textsc{direct} bounds what a
shadow-model attacker may realize after transfer. Shard Protection is
therefore a defensive insight, not a deployment guarantee: it limits
utility collapse but does not by itself address privacy. Future work
should study mechanism-aware selectors, partition privacy--robustness
tradeoffs, cross-architecture surrogate transfer, and privacy-oriented
attacks rather than only downstream-F1 attacks.
